%% ============================================================
%%  es/experience-es.tex  —  Experiencia Profesional
%% ============================================================

\section{Experiencia Profesional}

%% ── 1. IDM Technology / Scotiabank  (Marzo 2026 – Presente) ──
\cvEntryClient%
  {Desarrollador Fullstack Senior}%
  {IDM Technology}%
  {Scotiabank Peru}%
  {Plataforma de Pagos QR Dinámico}%
  {Marzo 2026 -- Presente \enspace\textbar\enspace Contrato 6 meses}

\begin{cvItems}
  \cvItem{Desarrollo de microservicios reactivos cloud-native con \textbf{Java 25, Spring Boot 4} y
    \textbf{Spring WebFlux} para una plataforma de pagos QR dinámico.}
  \cvItem{Diseño e implementación de microservicios para \textbf{generación de QR dinámico y estático}
    en pagos de billetera digital interoperables.}
  \cvItem{Construcción de \textbf{REST APIs} para gestión de QR, procesamiento de pagos y persistencia
    de transacciones.}
  \cvItem{Desarrollo de \textbf{microservicios de integración} entre servicios de Google Cloud y
    sistemas mainframe IBM AS/400 a través de DataQueues.}
  \cvItem{Implementación de mecanismos de autenticación y autorización bajo \textbf{estándares de
    seguridad bancaria} (Spring Security).}
  \cvItem{Diseño de capas de persistencia con \textbf{Microsoft SQL Server}; optimización de queries
    para cargas transaccionales de alto rendimiento.}
  \cvItem{Trabajo con \textbf{Google Cloud Platform (GCP)} y \textbf{Apigee} para gestión de APIs
    segura y basada en políticas.}
  \cvItem{Participación en pipelines de \textbf{Jenkins CI/CD} y mantenimiento de calidad de código
    con \textbf{SonarCloud}; resolución de vulnerabilidades identificadas en pruebas de penetración.}
  \cvItem{Colaboración en equipos Agile Scrum multifuncionales usando Jira, Confluence y Bitbucket.}
  \cvItem{Desarrollo de herramientas internas paralelas para el banco usando \textbf{React},
    \textbf{Electron}, \textbf{Node.js} y \textbf{Python} para soporte de flujos operativos y
    automatización interna.}
  \cvItem{Aplicación de \textbf{desarrollo asistido por IA} (Claude Code, GitHub Copilot) para
    acelerar scaffolding, revisión de código y entrega en todos los workstreams paralelos.}
\end{cvItems}

{\footnotesize\color{cvMuted}\textit{Stack:}~%
Java 25, Spring Boot 4, Spring WebFlux, Programación Reactiva, REST APIs, SQL Server,
IBM AS/400, DataQueues, Spring Security, GCP, Apigee, Jenkins, SonarCloud, Bitbucket,
React, Electron, Node.js, Python, Claude Code, GitHub Copilot.}

\vspace{4pt}

%% ── 2. DaCodes — AgentCloud  (Enero 2025 – Febrero 2026) ────
\cvEntryClient%
  {Ingeniero de Software Líder}%
  {DaCodes}%
  {HealthCare}%
  {AgentCloud}%
  {Enero 2025 -- Febrero 2026}

\begin{cvItems}
  \cvItem{Liderazgo en arquitectura y desarrollo full-stack de \textbf{AgentCloud}, plataforma
    empresarial de orquestación de agentes IA para clientes en Estados Unidos.}
  \cvItem{Diseño de microservicios cloud-native en \textbf{AWS} (Lambda, EKS, SQS) con Java y
    Spring Boot; contenedorizados con Docker y desplegados vía Kubernetes.}
  \cvItem{Construcción de frontend en \textbf{React} para configuración, monitoreo y despliegue
    en tiempo real de flujos de agentes IA para usuarios de negocio.}
  \cvItem{Definición de estándares de ingeniería, revisiones de código y liderazgo de equipo
    distribuido en ceremonias Agile diarias.}
  \cvItem{Implementación de \textbf{pipelines CI/CD} (GitHub Actions) y servicios backend
    \textbf{event-driven} con Apache Kafka para coordinación de agentes en tiempo real.}
  \cvItem{Uso de \textbf{agentes de codificación IA} (Claude Code, GitHub Copilot, Cursor) para
    scaffolding backend, revisiones de código automatizadas e ingeniería de prompts dentro de
    la plataforma de orquestación de agentes.}
\end{cvItems}

\vspace{4pt}

%% ── 3. ExSquared / SambaSafety — Qorta  (Marzo 2023 – Dic 2024)
\cvEntryClient%
  {Ingeniero Full Stack Senior}%
  {ExSquared}%
  {SambaSafety}%
  {Qorta}%
  {Marzo 2023 -- Diciembre 2024}

\begin{cvItems}
  \cvItem{Contribución a \textbf{Qorta}, plataforma de monitoreo de riesgo de conductores de
    SambaSafety usada por operadores de flotas en Norteamérica.}
  \cvItem{Construcción de \textbf{REST APIs} escalables con \textbf{NestJS / Node.js} e integración
    con proveedores de registros vehiculares (MVR) en más de 50 jurisdicciones de EE.UU.}
  \cvItem{Desarrollo de componentes de dashboard en \textbf{Angular} para visualización de puntajes
    de riesgo, reportes de cumplimiento y gestión de alertas en tiempo real.}
  \cvItem{Implementación de flujos de procesamiento de datos con \textbf{Kafka} y \textbf{Redis};
    optimización de queries \textbf{PostgreSQL} para grandes volúmenes de datos de conductores.}
  \cvItem{Mantenimiento de cobertura de pruebas superior al 85\% con \textbf{Jest} y \textbf{Karma};
    colaboración en equipos Agile Scrum distribuidos con ingenieros en EE.UU.}
\end{cvItems}

\vspace{4pt}

%% ── 4. Globant / Banco GNB Paraguay  (Junio 2022 – Mayo 2023) ──
\cvEntryClient%
  {Ingeniero Full Stack Senior}%
  {Globant}%
  {Banco GNB Paraguay}%
  {Plataforma de Banca Digital}%
  {Junio 2022 -- Mayo 2023}

\begin{cvItems}
  \cvItem{Entrega de funcionalidades de banca digital como ingeniero embebido de Globant en el
    equipo de producto principal de Banco GNB Paraguay.}
  \cvItem{Desarrollo de microservicios en \textbf{Java / Spring Boot} para gestión de cuentas,
    procesamiento de transacciones y onboarding de clientes.}
  \cvItem{Construcción de módulos \textbf{Angular} para el portal online del banco; integración con
    sistemas bancarios core mediante \textbf{REST APIs} y capas de mensajería.}
  \cvItem{Aplicación de prácticas de desarrollo orientadas a seguridad para cumplimiento regulatorio
    financiero en el sector bancario paraguayo.}
\end{cvItems}

%% ── Experiencia Adicional (condensada) ──────────────────────
\section{Experiencia Adicional}

{\small
\cvMiniEntry{Ingeniero Full Stack}{The Bridge Social (Remoto, EE.UU.)}
\cvMiniEntry{Ingeniero Full Stack}{INDRA\enspace\textbar\enspace Clientes: BCP, RIMAC}
\cvMiniEntry{Ingeniero Full Stack}{Michael Page\enspace\textbar\enspace Cliente: Intercorp}
\cvMiniEntry{Ingeniero Full Stack}{Zoluxiones\enspace\textbar\enspace Cliente: SURA}
\cvMiniEntry{Ingeniero Full Stack}{Experis\enspace\textbar\enspace Cliente: Equifax}
\cvMiniEntry{Ingeniero Full Stack}{Olva Courier}
\cvMiniEntry{Ingeniero Full Stack}{LimaW}
}

%% ── Logros Destacados ────────────────────────────────────────
\section{Logros Destacados}

\begin{cvItems}
  \cvItem{\textbf{Más de 12 años} de ingeniería de software profesional en Salud, Banca,
    Seguros, Retail y Logística.}
  \cvItem{Entrega de soluciones en producción para clientes de primer nivel: \textbf{Scotiabank},
    SambaSafety, Banco GNB Paraguay, Equifax, Intercorp, SURA, RIMAC y BCP.}
  \cvItem{Liderazgo de equipos de ingeniería distribuidos en proyectos 100\% remotos en
    \textbf{Estados Unidos, América Latina y el Caribe}.}
  \cvItem{Experiencia comprobada en \textbf{modernización de sistemas legados bancarios y de seguros}
    hacia arquitecturas de microservicios cloud-native.}
  \cvItem{Entrega consistente en entornos regulados de alta disponibilidad con estrictos requisitos
    de seguridad y cumplimiento normativo.}
\end{cvItems}
